%--------------------------------------------------
%  Topología – Clase 3
%--------------------------------------------------
\subsection{Bases y topologías generadas}

\paragraph{\textbf{Proposición.}}
Sea $\mathcal{B}$ una colección de subconjuntos de un conjunto $X$ tal que
\begin{enumerate}
  \item para todo $x\in X$ existe $B\in\mathcal{B}$ con $x\in B$;
  \item para cualesquiera $B_{1},B_{2}\in\mathcal{B}$ y todo $x\in B_{1}\cap B_{2}$
        existe $B_{3}\in\mathcal{B}$ con $x\in B_{3}\subseteq B_{1}\cap B_{2}$.
\end{enumerate}
Entonces $\mathcal{B}$ es una base de una topología en $X$.

\paragraph{\textbf{Ejemplo}}
En $\mathbb{R}^{n}$ con la topología usual, sea
\[
  \mathcal{B}=
  \Bigl\{
    \prod_{i=1}^{n}(a_{i},b_{i}) : a_{i},b_{i}\in\mathbb{R},\;i=1,\dots,n
  \Bigr\}.
\]
\begin{enumerate}
  \item Para $x=(x_{1},\dots,x_{n})$, tomamos
        \[
          B=\prod_{i=1}^{n}(x_{i}-1,x_{i}+1)\in\mathcal{B},
          \quad x\in B.
        \]
  \item Sean
        \[
          B_{1}=\prod_{i=1}^{n}(a_{i},b_{i}),
          \quad
          B_{2}=\prod_{i=1}^{n}(c_{i},d_{i}).
        \]
        Si $x\in B_{1}\cap B_{2}$, entonces
        \[
          B_{3}=\prod_{i=1}^{n}
                 \bigl(
                   \max\{a_{i},c_{i}\},
                   \min\{b_{i},d_{i}\}
                 \bigr)
          \;\in\mathcal{B},
          \quad x\in B_{3}\subseteq B_{1}\cap B_{2}.
        \]
\end{enumerate}

Por lo tanto $\mathcal{B}$ es una base para alguna topología en $\mathbb R^n$. 
\subsubsection{Relación con la topología usual}

\paragraph{\textbf{$\mathcal T_{\mathcal{B}}\subseteq\mathcal T_{u}$.}}

\textbf{Ejercicio.}

\paragraph{\textbf{$\mathcal T_{u}\subseteq\mathcal T_{\mathcal{B}}$.}}
Para cada $x\in\mathbb{R}^{n}$ y $\varepsilon>0$ sea
\[
  U=\prod_{i=1}^{n}\!\bigl(x_{i}-\tfrac{\varepsilon}{n},\,x_{i}+\tfrac{\varepsilon}{n}\bigr)
  \;\in\mathcal{B}.
\]
Si $y=(y_{1},\dots,y_{n})\in U$ entonces $|x_{i}-y_{i}|<\tfrac{\varepsilon}{n}$
para todo $i$ y
\[
  \|x-y\|_{1}
  =\sum_{i=1}^{n}|x_{i}-y_{i}|<\varepsilon,
\]
luego $U\subset B_{\varepsilon}(x) = \{y\in\mathbb{R}^{n}:\|x-y\|_{1}<\varepsilon\}$.

\paragraph{\textbf{Proposición.}}
Sea $(X,\mathcal T)$ un espacio topológico y $\mathcal{B}$ una base de $\mathcal T$. Entonces $\mathcal T=\mathcal T_{\mathcal{B}}$ si y solo si:

\begin{itemize}
  \item $\mathcal T_\mathcal{B} \subseteq \mathcal T$
  \item $\forall U \in \mathcal T,\forall x \in U, \exists B \in \mathcal B:  x \in B \subseteq U$ 
\end{itemize}

%--------------------------------------------------
\subsection{Vecindades y subespacios}

\paragraph{\textbf{Definición (Vecindad).}}
Sea $x\in X$. Un conjunto $A\subseteq X$ es una \emph{vecindad} de $x$
si $x\in A$ y existe $U\in\mathcal T$ con $x\in U\subseteq A$. Cuando además
$A\in\mathcal T$ se la llama \emph{vecindad abierta}.

\paragraph{\textbf{Ejercicio.}}
Un subconjunto $A\subseteq X$ es abierto \emph{ssi} es vecindad abierta de todos sus puntos.

\subsubsection{Topología de subespacio}

\paragraph{\textbf{Definición.}}
Para $(X,\mathcal T)$ espacio topológico y $A\subseteq X$, definimos
\[
  \mathcal T_{A} = \{\,U\subseteq A: U=A\cap V\text{ con }V\in\mathcal T\}.
\]

Veamos que $\mathcal T_A$ es una topología:

\begin{enumerate}
\item $\emptyset \in \mathcal T_A$ pues $\emptyset = A \cap \emptyset$ y $\emptyset \in \mathcal T_A$.
\item $A \in \mathcal T_A$ pues $A = A \cap X$ y $X \in \mathcal T$.
\item Sea $\{U_i\}_{i \in I}$ una familia de elementos de $\mathcal T_A$. Entonces
\[
U_i = A \cap V_i \quad \text{para cada } i \in I
\]
con $V_i \in \mathcal T$. Luego:
\[
\bigcup_{i \in I} U_i = \bigcup_{i \in I} (A \cap V_i) = A \cap \left(\bigcup_{i \in I} V_i\right)
\]
y como $\bigcup_{i \in I} V_i \in \mathcal T$, entonces $A \cap \left(\bigcup_{i \in I} V_i\right) \in \mathcal T_A$.
\item (Ejercicio) La intersección finita de elementos de $\mathcal T_A$ pertenece a $\mathcal T_A$.
\end{enumerate}

\paragraph{\textbf{Definición.}} Si $B\subseteq A$ tal que $B \in \tau_A$, decimos que $B$ es abierto (cerrado) en $A$  si $B (A-B) \in \tau_A$  .


\paragraph{\textbf{Ejemplo.}}
Para $A=\overline{B_{1}}(0)\subset\mathbb{R}^{n}$, el conjunto
$\overline{B_{1}}(0)$ es abierto en $A$ pero no en $\mathbb{R}^{n}$.

\paragraph{\textbf{Proposición.}}
Si $A$ es abierto en $X$ y $B$ es abierto en $A$, entonces $B$ es abierto en $X$.

\paragraph{\textbf{Demostración.}}
Dado que $B \in  \mathcal T_A$ entonces existe $V \in \mathcal T$ tal que $V \overset{ab} \subseteq X$, $B = A \cap V$, y como $A \in \mathcal T$, entonces $B = A \cap V \in \mathcal T$. Es decir existen dos abiertos $A$ y $V$ tales que $B = A \cap V$, y como la intersección de abiertos es abierto, entonces $B$ es abierto en $X$.

\paragraph{\textbf{Proposición.}}
Sea $(X,d)$ un espacio métrico y $A\subseteq X$. La topología $\mathcal T_{A}$
coincide con la inducida por la métrica restringida $d|_{A}$.

\paragraph{\textbf{Demostración.}}

\begin{itemize}
  \item \textbf{¿Por qué $\mathcal T_A \supseteq \mathcal T_{d|A}$?}  
  Se debe probar que para todo $\epsilon>0$ y para todo $x\in A$:
  $$
  B_\epsilon^{d|_A}(x) = \{ y \in A : d(x,y) < \epsilon \} = B_{\epsilon}^d(x) \cap A \in \mathcal T_A
  $$
  donde $B_\epsilon^d(x) = \{ y \in X : d(x,y) < \epsilon \}$.
  
  Luego, por definición, $B_\epsilon^d(x) \in \mathcal T$ y entonces $B_\epsilon^d(x) \cap A \in \mathcal T_A$.

  \item \textbf{¿Por qué $\mathcal T_{d|A} \supseteq \mathcal T_A$?}  
  Sea $U \in \mathcal T_A$, es decir, $U = V \cap A$ para algún $V \in \mathcal T$.  
  Para cada $x \in U$, existe $\epsilon > 0$ tal que:
  $$
  B_\epsilon^d(x) \subseteq V
  $$
  Luego:
  $$
  B_\epsilon^{d|A}(x) = B_\epsilon^d(x) \cap A \subseteq V \cap A = U
  $$
  por lo tanto, $U$ es abierto en $\mathcal T_{d|A}$.
\end{itemize}

De todo lo anterior concluimos que $\mathcal T_A = \mathcal T_{d|_A}$

%--------------------------------------------------
\subsubsection{Subespacios discretos}

\paragraph{\textbf{Definición.}}
Un subconjunto $A\subseteq X$ es \emph{discreto} si $\mathcal T_{A}$ es la topología discreta sobre $A$.

\paragraph{\textbf{Ejemplo.}}
\begin{enumerate}[label=(\alph*)]
  \item $\mathbb{Z}$,
  \item $\{y_{n}:n\in\mathbb{Z}^{+}\}$,
  \item $\{y_{n}:n\in\mathbb{Z}^{+}\}\cup\{0\}$,
  \item $\mathbb{Q}$.
\end{enumerate}

$$
\{n\} = \left(n-1,n+1\right) \cap \mathbb Z
$$
 entonces $\{n\} \subseteq^{ab} \mathbb Z$ para todo $n$ tiene la topologia discreta.


$$
\left\{\frac 1{n+1}\right\} = \left(\frac 1{n+2},\frac 1{n+1}\right) \cap A
$$
con $n\geq 1$ 


$$
\{
    1
\} =  \left(\frac 12,2\right) \cap A$$


para todo $(a,b)$ tq $0 \in (a,b)$ existe $n \in \mathbb Z^+$ tq $\frac 1n \in (a,b)$, entonces $\{0\}$ no es abierto en $B$. Entonces $B$ no tienen la topologia discreta,



%--------------------------------------------------
\subsection{Clausura en subespacios}

\paragraph{\textbf{Lema.}}
Sea $Y\subseteq X$ un subespacio y $A\subseteq Y$. Entonces
\[
  \overline{A}^{Y} = \overline{A}^{X}\cap Y.
\]

\paragraph{\textbf{Demostración.}}
\begin{itemize}
  \item [$\subseteq$] Sea $x$ un elemento en la clausura de $A$ en $Y$. Sea $U$ una vecindad abierta de $x$ en $Y$:

  \begin{align*}
  U \cap A \neq \emptyset
  \end{align*}
  
  ademas $U = V \cap Y$, por lo que:

  \begin{align*}
  U \cap A &= (V \cap Y) \cap A \\
  &= (V \cap A) \cap Y
  \end{align*}
  % <!--- con $V$ vecindad abierta de $x$ en $X$, luego 
  
  
  % 
  % \begin{align*}
  % V\cap A 
  % \end{align*}
  % > -->
  
  Sea $V$ vecindad abierta de $x$ en $X$. Dado que $V\cap A \neq \emptyset$, pues en otro caso:
  \begin{align*}
  U \cap A &= (V \cap Y) \cap A \\
  &= \emptyset \cap A = \emptyset
  \end{align*}

  Entonces $ x \in \overline{A}^{X}$, y como $x \in Y$, entonces $x \in \overline{A}^{X} \cap Y$.

  \item [$\supseteq$] $x$ un elemento en la clausura de $A$ en $X\cap Y$

  Sea $U$ una vecindad abierta de $x$ en $Y$, queremos mostrar que $U\cap A \neq \emptyset$ 

  $U = V \cap Y$ con $V$ vecindad abierta de $x$ en $X$. Sabemos que:
  \begin{align*}
  V \cap A \neq \emptyset
  \end{align*}
  entonces
  \begin{align*}
  U\cap A = (V\cap Y) \cap A  = V \cap A \neq \emptyset
  \end{align*}
  Luego $x$ es un elemento en la clausura de $A$ en $Y$.
\end{itemize}


%--------------------------------------------------
\subsection{Ejercicios}

(Viro--Ivanov~2008, pág.~17)\,: [3'4], [3'5], 4.A, 4.E, [4'12], [4'16x], [6'3], [6'9], [6'12]--sust.\ 2.Mx--, [2'13x].

%--------------------------------------------------
\subsection{Homeomorfismos}

\paragraph{\textbf{Definición.}}
Una función biyectiva $f\colon X\to Y$ entre espacios topológicos es un \emph{homeomorfismo} si $f$ es continua y $f^{-1}$ también lo es.